\documentclass[10pt,a4paper]{amsart}
\usepackage[margin=19mm]{geometry}
\usepackage{amsmath,amssymb,mathtools}
\usepackage[scr=rsfs,cal=euler]{mathalfa}
\usepackage{xfrac}
\usepackage[unicode,hidelinks,bookmarks=false]{hyperref}
\newcommand{\ie}{i.e.\ }
\newcommand{\cf}{cf.\ }
\newcommand{\resp}{respectively\ }
\newcommand{\Subsection}{\subsection}
\providecommand{\nobreakdash}{\nobreak\hbox{-}}
\setlength{\emergencystretch}{2em}
\multlinegap0pt
\usepackage{amssymb}
\usepackage{enumerate}
\usepackage{dsfont}
\newtheorem{remark}{Remark}
\newtheorem{theorem}{Theorem}
\newtheorem{definition}{Definition}
\newtheorem{lemma}{Lemma}
\newcommand{\Ba}{\mathscr B_z^{(a)}}
\newcommand{\C}{\mathbb{C}}
\def \R {\mathbb{R}}
\def \RN {\mathbb{R}_+^{d+1}}
\newcommand{\Rd}{\mathbb R^d}
\newcommand{\XT}{\mathcal{X}_T}
\newcommand{\Xs}{\mathcal X}
\newcommand{\p}{\partial}
\newcommand{\ve}{\varepsilon}
\title{Strichartz estimates for Schr\"odinger equations with nonlinear boundary interactions}
\author{Nicola Garofalo, Gigliola Staffilani}
\date{Source: arXiv:2605.23229v2 (CC BY 4.0)}
\begin{document}
\maketitle

\noindent\emph{Source and licence.} Strichartz estimates for Schr\"odinger equations with nonlinear boundary interactions. Version arXiv:2605.23229v2 (CC BY 4.0), distributed by arXiv under the Creative Commons Attribution 4.0 International licence, http://creativecommons.org/licenses/by/4.0/, as recorded in the arXiv licence metadata for this exact version and shown on the abstract page https://arxiv.org/abs/2605.23229v2. Full licence notice: \texttt{licenses/arxiv-2605.23229-CC-BY-4.0.txt}.\par\smallskip

\noindent\textit{Editorial scope.} Counted unit: the article's theorem T:well (source0.tex lines 690--703). The article's complete original proof of it is delivered with it (source0.tex lines 2601--2856, one \texttt{proof} environment of 256 lines, which the article places away from the statement). No mathematics is written, repaired or re-derived: the statement and the proof are the article's own lines, in the article's own notation, and no retained line is substituted or re-flowed.  Retained uncounted as the source-local context the proof needs: lines 359--365; lines 383--389; lines 489--491; lines 500--502; lines 555--558; lines 592--601; lines 613--623; lines 629--643; lines 652--660; lines 654--656; lines 668--670; lines 672--674; lines 680--686; lines 709--721; lines 913--915; lines 1695--1697; lines 1699--1701; lines 2066--2081; lines 2221--2227; lines 2358--2364; lines 2439--2439; lines 2493--2495; lines 2498--2498; lines 2545--2547; lines 2564--2574; lines 2880--2882.  Nothing was omitted from any retained span, so the package carries no omission record.  No retained line contains a citation, so the wrapper carries no bibliography and no bibliography entry.  No retained line contains a figure environment, an image-inclusion command or drawing code, so no image is delivered and the package declares no asset dependency.  Editorial formatting only, in addition: an AMS \texttt{amsart} preamble that restates the article's own declarations and macros used by the retained lines, namely the environments \texttt{remark}, \texttt{theorem}, \texttt{definition}, \texttt{lemma} on separate counters, together with the standard packages those lines need.  The retained environments print in this document as Remark 1, Theorem 1, Definition 1, Lemma 1 in order of appearance, whereas the article prints the same items with the numbering of its own full document.  The complete native source of the article as distributed in the arXiv e-print for this exact version is bundled unchanged as \texttt{sources/201197/source0.tex}.
\begin{equation}\label{cp02}
\begin{cases}
\p_t U - i(\Delta_x U + \Ba U) = F(X,t), & \ \ X=(x,z)\in \mathbb{R}^{d+1}_+,\ t\in \mathbb R, \\
\displaystyle \lim_{z\to 0^+} z^a \p_z U(X,t) = -\mu |U(x,0,t)|^{p-1}U(x,0,t), & \ \ x\in \Rd,\ t\in \R, \\
U(X,0)=u_0(X), 
\end{cases}
\end{equation}

\begin{equation}\label{nozeroin}
\begin{cases}
\partial_tU-i(\Delta_x U+\Ba U)=F(X,t),\\
\displaystyle \lim_{z\to0^+} z^a\partial_zU(X,t)=\Phi(x,t),\\
U(X,0)=u_0(X).
\end{cases}
\end{equation}

\begin{equation}\label{U0}
U(X,t) = \mathbb T^\star_a(u_0)(X,t) + \mathbb D_a(F)(X,t) + \Theta^\star_a(\Phi)(X,t),
\end{equation}

\begin{equation}\label{Thetastar}
\Theta_a^\star(\Phi)(X,t) := -i \int_0^{t} \int_{\Rd}  \mathbb S_a(X,Y_0,t-\tau) \Phi(y,\tau) dy d\tau.
\end{equation}

\begin{equation}\label{acinfty}
a\ge 0:\ \ \ \ \ \frac 2q +\frac dr = \frac{d+a+1}2;\ \ \ \ \ \ \  -1<a<0:\ \ \ \ \ 
\frac 2q +\frac dr = \frac{d+1}2.
\end{equation}

\begin{remark}\label{R:rr}
For future use, it will be important to keep in mind that the triple $(r,r,\infty)$ is always admissible, and we have: 
\begin{equation}\label{same}
q = r = \frac{2(d+2)}{d+a+1},\ \ \ \ \ \ \ \ \ \ \ q' = r' = \frac{2(d+2)}{d+3-a},\ \ \ \ \text{when}\ \ a\ge 0,
\end{equation}
and
\begin{equation}\label{same0}
q = r = \frac{2(d+2)}{d+1},\ \ \ \ \ \ \ \ \ \ \ q' = r' = \frac{2(d+2)}{d+3},\ \ \ \ \text{when}\ \ -1<a< 0.
\end{equation}
\end{remark}

\begin{theorem}\label{T:main}
Let $0\le a<1$ and suppose that the triple $(q,r,\infty)$ is admissible, i.e., that the first equation in \eqref{acinfty} holds. If $d+a-1>0$ we require that $r<\frac{2d}{d+a-1}$, if instead $d+a-1=0$, we assume that $(q,r,\infty)\not= (2,\infty,\infty)$.
 Then there exists a constant $C = C(d,a,r)>0$ such that the following Strichartz estimates hold for the solution $U$ to the problem \eqref{nozeroin}:
\begin{align}\label{Umain}
||U||_{L^\infty_t L^2_a(\RN)} \le C\ \left[||u_0||_{L^2_a(\RN)} + ||F||_{L^1_t L^2_{a}(\RN)} + ||\Phi||_{L^{q'}_t L^{r'}_x} \right],
\end{align}
and also
\begin{align}\label{Umain2}
||U||_{L^\infty_z L^q_t L^r_x } \le C\ \left[||u_0||_{L^2_a(\RN)} +  ||F||_{L^1_{a,z} L^{q'}_t L^{r'}_x } + ||\Phi||_{L^{q'}_t L^{r'}_x} \right].
\end{align}
\end{theorem}

\begin{theorem}\label{T:main2}
Let $-1<a<0$. Assume that $(q_\infty,r,\infty)$ and $(q,r,\infty)$ respectively satisfy the first and second equations in \eqref{acinfty}, i.e., 
\begin{equation}\label{uffaa}
\frac{2}{q_\infty} + \frac dr = \frac{d+a+1}2,\ \ \ \ \text{and}\ \ \ \ \ \frac 2{q} + \frac dr = \frac{d+1}2. 
\end{equation} 
Moreover, if $d>1$ we require that $2\le r < \frac{2d}{d-1}$, whereas when $d = 1$ we assume that $(q,r,\infty)\not= (2,\infty,\infty)$.
Then there exists a constant $C = C(d,a,r)>0$ such that the following Strichartz estimates hold for the solution $U$ to the problem \eqref{nozeroin}:
\begin{align}\label{Umain20}
||U||_{L^\infty_t L^2_a(\RN)} \le C\ \left[||u_0||_{L^2_a(\RN)} + ||F||_{L^1_t L^2_{a}(\RN)} + ||\Phi||_{L^{q'_\infty}_t L^{r'}_x} \right],
\end{align}
and also
\begin{align}\label{Umain22}
||U\ k||_{L^\infty_z L^{q+q_\infty}_t L^r_x } \le C\ \left[||u_0||_{L^2_a(\RN)} +  ||F \ k^{-1}||_{L^1_{a,z} L^{q'\cap q'_\infty}_t L^{r'}_x } + ||\Phi||_{L^{q'_\infty}_t L^{r'}_x} \right].
\end{align}
\end{theorem}

\begin{definition}\label{solution}
Let $\mu\in \C\setminus\{0\}$. Assume that $0\le a<1$ and that $(q,r,\infty)$ is admissible, with $(q,r,\infty)\not= (2,\infty,\infty)$ when $d=1$ and $a = 0$. Given $u_0\in L^2_a(\RN)$ and $F\in L^1_t L^2_a(\RN) \cap L^1_{a,z} L_t^{q'} L^{r'}_x$, we say that a function $U$ is a \emph{mild solution} to the Cauchy problem \eqref{cp02} if $U\in C(\R,L^2_a(\RN))$, $z\to U(\cdot,z,\cdot)$ belongs to $C^b_z L^q_tL^r_x$, and it satisfies the identity 
\begin{align}\label{sol1}
U(X,t) & = \mathbb T^\star_a(u_0)(X,t) +  \mathbb D_a(F)(X,t) - \mu\, \Theta^\star_a(|U|^{p-1} U(\cdot,0,\cdot))(X,t).
\end{align}
Note that, in view of \eqref{Thetastar} and the Neumann property of the representation \eqref{U0}, the identity \eqref{sol1} encodes precisely the boundary condition $\lim_{z\to 0^+} z^a\p_z U = -\mu|U|^{p-1}U$ in \eqref{cp02}.
If instead $-1<a<0$, given $0<T<\infty$, $u_0\in L^2_a(\RN)$ and $F\in L^1_T L^2_a(\RN)$ such that 
$F  k^{-1}\in L^1_{a,z} L^{q'\cap q'_\infty}_T L^{r'}_x $, we say that $U$ is a \emph{mild solution} on $[0,T]$ if $U\in C([0,T],L^2_a(\RN))$, $z \to U(\cdot,z,\cdot)k(z)$ belongs to $C^b_z L^q_T L^r_x$, and $U$ satisfies the identity \eqref{sol1}.
\end{definition} 

\begin{align}
U(X,t) & = \mathbb T^\star_a(u_0)(X,t) +  \mathbb D_a(F)(X,t) - \mu\, \Theta^\star_a(|U|^{p-1} U(\cdot,0,\cdot))(X,t).
\end{align}

\begin{equation}\label{pcpos}
p = p_c := 1 + \frac{2(1-a)}{d+a+1},\ \ \ \text{when}\ \ \ 0\le a<1,  
\end{equation}

\begin{equation}\label{pcneg}
p = p_c := 1 + \frac{2}{d+1},\ \ \ \ \ \text{when}\ \ \ -1<a<0.
\end{equation}

\begin{remark}\label{R:pc}
When the triple $(r,r,\infty)$ is admissible, the exponent $q=r$ is uniquely identified by the equations \eqref{same} or \eqref{same0}. One verifies that in either case we have $q=r = p_c+1$, the critical exponent in \eqref{pcpos}, \eqref{pcneg}. We also note that, if we set $\tilde q := p_c q'$ and $\tilde r := p_c r' $, we have (see also Lemma \ref{L:mass})
\[
\tilde q = \tilde r = p_c+1 = q = r.
\]
This will be important in the proof of $L^2_a$-mass critical case in Theorems \ref{T:well} and \ref{T:nega}.
\end{remark}

\begin{theorem}\label{T:nega}
Let $-1<a<0$ and, given $0<T<\infty$, $\mu\in \C$ and $1<p\le p_c$, consider the problem \eqref{cp02} on $\RN\times[0,T]$. Let $r = p+1$ and suppose that the triples $(q_\infty,r,\infty)$ and $(q,r,\infty)$ satisfy the equations \eqref{uffaa}, with $q>2$. Let $u_0\in L^2_a(\RN)$ and $F\in L^1_T L^2_a(\RN)$ be such that 
$F  k^{-1}\in L^1_{a,z} L^{q'\cap q'_\infty}_T L^{r'}_x $.
\begin{itemize}
\item[(1)] \underline{\emph{The $L^2_a$-mass critical case}}. 
If $p=p_c$, there exists $\ve_0 = \ve_0(d,a,\mu,T)>0$ such that, if   
\[
 C \max\{1,T^{\frac{1}{q}-\frac{1}{q_\infty}}\} \left[\|u_0\|_{L^2_a(\RN)} + ||F||_{L^1_T L^2_a(\RN)} + (1+T^{\frac{1}{q}-\frac{1}{q_\infty}}) ||F k^{-1}||_{L^1_{a,z} L^{q'\cap q'_\infty}_T L^{r'}_x}\right]\le \varepsilon_0,
\]
then there exists a mild solution $U\in C([0,T],L^2_a(\RN))$ to the problem \eqref{cp02}, such that $U k\in C^b_z L^q_T L^r_x$, unique in the ball $\{||U||_{\XT}\le 2\ve_0\}$ of the space $\XT$ in \eqref{Xneg}.
\item[(2)] \underline{\emph{The $L^2_a$-mass subcritical case}}. If $1<p<p_c$, then for any $\mu\in \C$ it is possible to select a small $0<T_0\le T$, depending only on $d, a, p, |\mu|, \|u_0\|_{L^2_a(\RN)}$ and the norms of $F$, for which there exists a mild solution $U\in C([0,T_0],L^2_a(\RN))$ to the problem \eqref{cp02} in $\RN\times [0,T_0]$, such that $U k\in C^b_z L^q_{T_0} L^r_x$, unique in $\mathcal X_{T_0}$.
\end{itemize}
\end{theorem}

\begin{equation}\label{Cz}
C^b_z L^q_t L^r_x =  C_z L^q_t L^r_x \cap L^\infty_z L^q_t L^r_x
\end{equation}

\begin{equation}\label{bulk00}
||\mathbb T^\star_a(u_0)||_{L^\infty_z L^q_t L^r_x } \le C(d,a,r)\ ||u_0||_{L^2_a(\RN)},
\end{equation}

\begin{equation}\label{Dabound}
||\mathbb D_a(F)||_{L^\infty_z L^q_t L^r_x } \le C(d,a,r)\  ||F||_{L^1_{a,z} L^{q'}_t L^{r'}_x }.
\end{equation}

\begin{theorem}\label{T:V}
Let $-1<a<1$ and suppose that $q, r$
satisfy
\begin{equation}\label{bdryad}
\frac 2q + \frac dr = \frac{d+a+1}2,
\end{equation}
with $r\ge 2$ and 
 $q>2$. Then there exists $C(d,a,r)>0$ such that for any $\Phi\in \mathscr S(\R^{d+1})$ one has
\begin{equation}\label{ok}
||\Theta_a^\star(\Phi)||_{L^\infty_t L^2_a(\RN)} \le C(d,a,r)\ ||\Phi||_{L^{q'}_t L^{r'}_x}.
\end{equation}
Furthermore, we also have
\begin{equation}\label{okk}
||\Theta_a^\star(\Phi)||_{L^\infty_{z} L^q_t L^r_x } \le C(d,a,r)\ ||\Phi||_{L^{q'}_t L^{r'}_x}.
\end{equation}
\end{theorem}

\begin{theorem}[Existence of traces]\label{T:gap}
Let $-1<a<1$ and suppose that $q, r$ satisfy \eqref{bdryad}, with $r\ge 2$ and $q>2$. Given $\Phi\in L^{q'}_t L^{r'}_x$,  
the function $\Theta^\star_a(\Phi)$ belongs to the space $C^b_z L^q_t L^r_x$. By this we mean that for every $z_0\ge 0$
\begin{equation}\label{cont}
\underset{z\to z_0}{\lim}\ ||\Theta^\star_a(\Phi)(\cdot,z,\cdot) - \Theta^\star_a(\Phi)(\cdot,z_0,\cdot)||_{L^q_t L^r_x} = 0.
\end{equation}
\end{theorem}

\begin{theorem}[Traces of the bulk terms]\label{T:gapbulk}
Let $0\le a<1$. Then for every $u_0\in L^2_a(\RN)$ and every $F$ with $||F||_{L^1_{a,z}L^{q'}_tL^{r'}_x}<\infty$ we have
\[
\mathbb T^\star_a(u_0),\ \mathbb D_a(F)\ \in\ C^b_zL^q_tL^r_x.
\]
If instead $-1<a<0$, then for every $u_0\in L^2_a(\RN)$ and every $F$ with $||Fk^{-1}||_{L^1_{a,z}L^{q'\cap q'_\infty}_tL^{r'}_x}<\infty$, the maps $z\to \mathbb T^\star_a(u_0)(\cdot,z,\cdot)k(z)$ and $z\to \mathbb D_a(F)(\cdot,z,\cdot)k(z)$ belong to $C^b_zL^{q+q_\infty}_tL^r_x$.
\end{theorem}

\section{Proof of Theorems \ref{T:main} and \ref{T:main2}}\label{S:main}

\begin{remark}\label{R:density}
Theorems \ref{T:main}, \ref{T:main2} and \ref{T:V} are stated for Schwartz data, but all the estimates are stable under limits. Consequently, the operators $\mathbb T^\star_a$, $\mathbb D_a$ and $\Theta^\star_a$ extend uniquely, by density, to bounded operators between the spaces appearing in the respective estimates, and for $u_0\in L^2_a(\RN)$, $\Phi\in L^{q'}_tL^{r'}_x$ and $F$ in the corresponding mixed-norm classes, the function $U$ defined by \eqref{U0} -- which is what we mean by the \emph{mild solution} of \eqref{nozeroin} with such data -- satisfies \eqref{Umain}, \eqref{Umain2}, \eqref{Umain20} and \eqref{Umain22} verbatim, as well as the trace properties of Theorems \ref{T:gap} and \ref{T:gapbulk}. In Section \ref{S:well} these extensions are used without further comment.
\end{remark}

\section{Well-posedness}\label{S:well}

\begin{lemma}\label{L:mass}
Let $a>-1$ and suppose that \eqref{cp02} is $L^2_a$-mass critical, with $p = p_c$. If the triple $(q,r,\infty)$ is admissible, and we let $\tilde q = p q'$, $\tilde r = p r'$, then also $(\tilde q,\tilde r,\infty)$ is admissible. 
\end{lemma}

\begin{lemma}[Conservation of mass]\label{L:masscon}
Let $0\le a<1$ and let $q, r$ satisfy \eqref{bdryad} with $r\ge2$, $q>2$; we recall that in this range \eqref{bdryad} coincides with the admissibility relation \eqref{acinfty}, so that Theorem \ref{T:gapbulk} applies with the same pair $(q,r)$. (We do not state the lemma in the anomalous range $-1<a<0$, where the trace of the bulk term provided by Theorem \ref{T:gapbulk} lives in the sum space $L^{q+q_\infty}_t L^r_x$ and the pairing below would require a different functional setting; the lemma is used only in the proof of Theorem \ref{T:well}, where $a\ge0$.) Given $u_0\in L^2_a(\RN)$ and $\Phi\in L^{q'}_tL^{r'}_x$, let $U = \mathbb T^\star_a(u_0) + \Theta^\star_a(\Phi)$ be the mild solution of the linear problem \eqref{nozeroin} with $F=0$ and Neumann datum $\Phi$ (see Remark \ref{R:density}), and let $u = U(\cdot,0,\cdot)\in L^q_tL^r_x$ be its boundary trace, which exists by Theorems \ref{T:gap} and \ref{T:gapbulk}. Then for every $t>0$
\begin{equation}\label{timeder}
\int_{\RN} |U(X,t)|^2 d\omega_a(X) = \int_{\RN} |u_0(X)|^2 d\omega_a(X)\ +\ 2\int_0^t\int_{\Rd} \Im\left(\overline{u(x,s)}\, \Phi(x,s)\right) dx\, ds .
\end{equation}
In particular, if $U$ is a mild solution of \eqref{cp02} with $F=0$, so that $\Phi = -\mu\,|u|^{p-1}u$, then
\begin{equation}\label{timeder2}
\int_{\RN} |U(X,t)|^2 d\omega_a(X) = \int_{\RN} |u_0(X)|^2 d\omega_a(X)\ -\ 2\,\Im(\mu) \int_0^t\int_{\Rd} |u(x,s)|^{p+1}\, dx\, ds,
\end{equation}
and if $\Im(\mu) = 0$ the \emph{mass} is constant in time.
\end{lemma}

\begin{equation}\label{Xneg}
\XT = \{U\in C([0,T],L^2_a(\RN))\ \text{such that}\  z\to U(\cdot,z,\cdot) k(z)\ \in C^b_z L^q_T L^r_x \},
\end{equation}

\begin{theorem}\label{T:well}
Let $0\le a<1$ and, given $\mu\in \C$ and $1<p\le p_c$, consider the problem \eqref{cp02}. Assume that $r = p+1$, and the triple $(q,r,\infty)$ is admissible with $q>2$. Suppose that $u_0\in L^2_a(\RN)$ and that  $F\in L^1_t L^2_a(\RN)\cap L^1_{a,z} L^{q'}_t L^{r'}_x $.
\begin{itemize}
\item[(1)] \underline{\emph{The $L^2_a$-mass critical case}}. If $p = p_c$ in \eqref{pcpos}, then with $q = r = p_c+1$ given by \eqref{same} there exists $\ve_0 = \ve_0(d,a,\mu)>0$ such that, if
\[
||u_0||_{L^2_a(\RN)} + ||F||_{L^1_t L^2_a(\RN)} + ||F||_{L^1_{a,z} L^{r'}_t L^{r'}_x} \le \ve_0,
\]
then there exists a global in time mild solution 
$U\in C(\R,L^2_a(\RN)) \cap C^{b}_z L^r_t L^r_x $ 
of the problem \eqref{cp02}, unique in the ball $\{||U||_{\Xs}\le 2\ve_0\}$ of the space $\Xs$ in \eqref{space}.
\item[(2)] \underline{\emph{The $L^2_a$-mass subcritical case}}. If $1<p<p_c$, then for any $\mu\in \C$ there exists a small $T = T(d,a,p,|\mu|,||u_0||_{L^2_a(\RN)},F)>0$ such that the problem \eqref{cp02} admits a mild solution $U\in C([0,T],L^2_a(\RN)) \cap C^{b}_z L^q_T L^r_x$, unique in $\XT$. 
If $F= 0$ and $\Im(\mu) = 0$, such $U$ can be extended to a mild solution on $[-M,M]$ for every $M>0$.
\end{itemize}
\end{theorem}

\begin{proof}[Proof of Theorem \ref{T:well}]
Assume that $a\ge 0$ and let $1<p\le p_c$, see \eqref{pcpos}. Suppose that $r = p+1$ and that $(q,r,\infty)$ satisfy the former equation in \eqref{acinfty}, with $q>2$. It is important to observe right away that, since $r-1 = p$, we must have $p r' = r$.

According to Theorems \ref{T:main} and \ref{T:gap}, and keeping \eqref{Cz} in mind, the appropriate Banach space for the problem is 
\begin{equation}\label{space}
\mathcal X= C(\R,L^2_a(\RN)) \cap C^{b}_z L^q_t L^r_x ,
\end{equation}
endowed with its natural norm 
\begin{equation}\label{norm}
||U||_{\Xs}= ||U||_{L^\infty_t L^2_a(\RN)} + ||U||_{L^\infty_z L^q_tL^r_x }.
\end{equation}
The couple $(\Xs,||\cdot||_{\Xs})$ is a complete Banach space. Indeed, each of the two spaces intersected in \eqref{space} is complete. For $C(\R,L^2_a(\RN))$, endowed with $||\cdot||_{L^\infty_t L^2_a}$, this is the standard fact that a uniform limit of continuous Banach space valued maps is continuous. The same argument in the variable $z$ applies to $C^b_z L^q_t L^r_x$, and shows that it is a closed subspace of $L^\infty_z L^q_t L^r_x$: here it is essential that, according to \eqref{Cz}, its norm be the genuine supremum over $z\ge 0$, and not an essential supremum. Finally, both spaces are continuously embedded in $L^1_{\operatorname{loc}}(\RN\times \R)$, and therefore a sequence which is Cauchy for \eqref{norm} has one and the same limit in the two of them. Consequently, $\Xs$ endowed with \eqref{norm} is complete.
In the sequel, given a function $U\in \Xs$, we denote for brevity  $|U|^{p-1} U(0) = |U(\cdot,0,\cdot)|^{p-1} U(\cdot,0,\cdot)$. Since $pr' = r$, we have
\begin{equation}\label{UU}
|||U|^{p-1} U(0)||_{L^{r'}_x} = ||U(0)||_{L^r_x}^p.
\end{equation}
Furthermore, if $U, V\in \Xs$, then we use the well-known inequality
\[
|\ |v|^{p-1} v - |w|^{p-1} w\ |\le C(p) \left[|v|^{p-1} + |w|^{p-1}\right] |v - w|,\ \ \ \ \ \ v, w\in \C,
\]
and H\"older inequality, to bound
\begin{align}\label{UUU}
|| |U|^{p-1} U(0) - |V|^{p-1} V(0)||_{L^{r'}_x} & \le C(p) \big\||U(0)| + |V(0)|\big\|^{p-1}_{L^{pr'}_x} ||U(0)-V(0)||_{L^{pr'}_x}
\\
& = C(p) \big\||U(0)| + |V(0)|\big\|^{p-1}_{L^{r}_x}  ||U(0)-V(0)||_{L^{r}_x}.
\notag
\end{align} 
 



In view of \eqref{sol1} in Definition \ref{solution}, given $u_0\in L^2_a(\RN)$ and $F\in L^1_t L^2_a(\RN) \cap L^1_{a,z} L_t^{q'} L^{r'}_x$, we now consider the operator: 
\begin{align}\label{La}
\Lambda_a(U)(X,t) & := \mathbb T_a^\star(u_0)(X,t) + \mathbb D_a(F)(X,t) - \mu\, \Theta_a^\star(|U|^{p-1} U(0))(X,t).
\end{align}
Then a function $U\in \Xs$ is a mild solution of the problem \eqref{cp02} if it is a fixed point of \eqref{La}, i.e., we have for every $(X,t)\in \RN\times \R$, with $t\not= 0$:
\begin{equation}\label{LaU}
U(X,t) = \Lambda_a(U)(X,t).
\end{equation}
From \eqref{La}, arguing as in the proof of Theorem \ref{T:main} in Section \ref{S:main}, we obtain the estimate  
\begin{align}\label{ini}
& ||\Lambda_a(U)||_{L^\infty_t L^2_a(\RN)} \le ||u_0||_{L^2_a(\RN)} + ||F||_{L^1_t L^2_a(\RN)}  + |\mu|\ ||\Theta^\star_a(|U|^{p-1}U(0))||_{L^\infty_t L^2_a(\RN)}
\\
& \le ||u_0||_{L^2_a(\RN)} + ||F||_{L^1_t L^2_a(\RN)}  + C_1\ |\mu|\ |||U|^{p-1}U(0)||_{L^{q'}_t L^{r'}_x},
\notag
\end{align}
where in the second inequality we have used \eqref{ok} in Theorem \ref{T:V}, and we have let $C_1 = C_1(d,a,r)>0$. 

Furthermore, by \eqref{bulk00}  and \eqref{Dabound}, we have for some $C_2 = C_2(d,a,r)>0$,
\begin{align}\label{bulkotto}
||\mathbb T^\star_a(u_0)||_{L^\infty_z L^q_t L^r_x} + ||\mathbb D_a(F)||_{L^\infty_z L^q_t L^r_x}\le C_2 \left[||u_0||_{L^2_a(\RN)} + ||F||_{L^1_{a,z} L^{q'}_t L^{r'}_x}\right].
\end{align}
From \eqref{norm}, \eqref{La}, \eqref{ini} and \eqref{bulkotto} we find for some $C_3 = C_3(d,a,r)>0$
\begin{align}\label{stelle1}
||\Lambda_a(U)||_{\Xs} & = ||\Lambda_a(U)||_{L^\infty_t L^2_a(\RN)} + ||\Lambda_a(U)||_{L^\infty_z L^q_t L^r_x}
\\
& \le C_3 \left[||u_0||_{L^2_a(\RN)} + ||F||_{L^1_t L^2_a(\RN)} + ||F||_{L^1_{a,z} L^{q'}_t L^{r'}_x}\right]
\notag
\\
& +  C_1\ |\mu|\ |||U|^{p-1}U(0)||_{L^{q'}_t L^{r'}_x}.
\notag
\end{align}
At this point we divide the discussion into the two cases (1) and (2). 

\medskip 

\noindent \underline{(1) The $L^2_a$-mass critical case.} We begin by observing that, if $p = p_c$ in \eqref{pcpos}, then 
\[
r = p+1 = \frac{2(d+2)}{d+a+1},
\]
and the compatibility assumption \eqref{acinfty} of the triple $(q,r,\infty)$ gives
\[
q = \frac{2(d+2)}{d+a+1} = r,
\]
see Remarks \ref{R:rr} and \ref{R:pc}.
Therefore, we have $r = r(d,a)$ and also $pq' = p r' = q$. From this fact and \eqref{UU} we thus find
\begin{align}\label{ini2}
|||U|^{p-1}U(0)||_{L^{q'}_t L^{r'}_x} = ||U(\cdot,0,\cdot)||^p_{L^{q}_t L^{r}_x} \le ||U||^p_{L^\infty_z L^{q}_t L^{r}_x }\le ||U||_{\Xs}^p,
\end{align}
where in the second inequality we have used the trivial bound
\begin{equation}\label{triv}
||U(\cdot,0,\cdot)||_{L^{q}_t L^{r}_x} \le ||U||_{L^\infty_z L^{q}_t L^{r}_x }.
\end{equation} 
Inserting \eqref{ini2} into \eqref{stelle1}, we reach the conclusion that 
\begin{align}\label{stelle5}
&  ||\Lambda_a(U)||_{\Xs} \le C_3 \left[||u_0||_{L^2_a(\RN)} + ||F||_{L^1_t L^2_a(\RN)} + ||F||_{L^1_{a,z} L^{q'}_t L^{r'}_x }\right] + C_1\ |\mu|\ ||U||_{\Xs}^p.
\end{align}


Our goal is to prove that there exists $\rho = \rho(d,a,\mu)>0$ small enough, such that if 
\[
B_\rho = \{U\in \Xs\mid ||U||_{\Xs}\le \rho\},
\]
then:
\begin{itemize}
\item[(i)] $\Lambda_a : B_\rho\ \longrightarrow\ B_\rho$;
\\
\item[(ii)] $||\Lambda_a(U) -\Lambda_a(V)||_{\Xs} \le \frac 12\ ||U-V||_{\Xs}$,
\end{itemize}
i.e. $\Lambda_a$ is a contraction. 

\vskip 0.2in

\noindent \underline{Proof of (i):} Suppose we choose $\rho = \rho(d,a,\mu)>0$ such that
\begin{equation}\label{rho}
C_1\ |\mu|\ \rho^p \le \frac{\rho}2\ \Longleftrightarrow\ 0< \rho \le (2 C_1\ |\mu|)^{-\frac{1}{p-1}}.  
\end{equation}
If $U\in B_\rho$, we thus have 
\[
C_1\ |\mu|\ ||U||_{\Xs}^p \le C_1\ |\mu|\ \rho^p \le \frac{\rho}2.
\]
It is thus clear from \eqref{stelle5} that, if we set $\ve_0 := \frac{\rho}2$, then 
\[
||U||_{\Xs}\le \rho\ \Longrightarrow\ ||\Lambda_a(U)||_{\Xs}\le \rho,
\]
provided that
\begin{equation}\label{smallid}
C_3 \left[||u_0||_{L^2_a(\RN)} + ||F||_{L^1_t L^2_a(\RN)} + ||F||_{L^1_{a,z} L^{q'}_t L^{r'}_x }\right]\  \le\ \ve_0.
\end{equation}
We can thus achieve (i) if \eqref{smallid} is true. We now turn to the 

\vskip 0.2in

\noindent \underline{Proof of (ii):} From the definition \eqref{La} we have
\begin{align}\label{Larap}
\Lambda_a(U) - \Lambda_a(V) = -\mu\, \Theta_a^\star\left(|U|^{p-1} U(0) - |V|^{p-1} V(0)\right).
\end{align}
From the boundary Strichartz estimate \eqref{okk}, we find
\begin{equation}\label{okkk}
||\Theta_a^\star\left(|U|^{p-1} U - |V|^{p-1} V\right)||_{L^\infty_{z} L^q_t L^r_x } \le C\ |||U|^{p-1} U(0) - |V|^{p-1} V(0)||_{L^{q'}_t L^{r'}_x}.
\end{equation}
Using \eqref{UUU}, H\"older inequality, and the fact that $pq' = q$, we easily obtain
\begin{align}\label{U4}
 || |U|^{p-1} U(0) - |V|^{p-1} V(0)||_{L^{q'}_t L^{r'}_x} & \le  C(p) \big\||U(0)| + |V(0)|\big\|^{p-1}_{L^q_t L^{r}_x}  ||U(0)-V(0)||_{L^q_t L^{r}_x}
 \\
& \le C(p) \left[||U(0)||^{p-1}_{L^q_t L^{r}_x}  + ||V(0)||^{p-1}_{L^q_t L^{r}_x}\right]  ||U(0)-V(0)||_{L^q_t L^{r}_x}.
\notag
\end{align}
From \eqref{triv}, \eqref{Larap}, \eqref{okkk} and \eqref{U4}, we conclude for $C_5 = C_5(d,a)>0$,
\begin{equation}\label{okkkk}
||\Lambda_a(U) - \Lambda_a(V)||_{L^\infty_{z} L^q_t L^r_x } \le C_5\ |\mu|\ \left[||U||_{L^\infty_z L^{q}_t L^{r}_x } + ||V||_{L^\infty_z L^{q}_t L^{r}_x}\right]^{p-1} ||U-V||_{L^\infty_z L^{q}_t L^{r}_x},
\end{equation}

If instead we invoke \eqref{ok}, from \eqref{Larap} and proceeding exactly as above, we find
\begin{align}\label{stelle6}
||\Lambda_a(U) - \Lambda_a(V)||_{L^\infty_t L^2_a(\RN)} \le C_6\ |\mu| \left[||U||_{ L^\infty_z L^{q}_t L^{r}_x} + ||V||_{ L^\infty_z L^{q}_t L^{r}_x}\right]^{p-1} ||U-V||_{ L^\infty_z L^{q}_t L^{r}_x}. 
\end{align}
From \eqref{norm}, \eqref{okkkk} and \eqref{stelle6} we conclude that, if $U, V\in B_\rho$, then
\[
||\Lambda_a(U) - \Lambda_a(V)||_{\Xs} \le C_7 |\mu| (2\rho)^{p-1} ||U - V||_{\Xs},
\]
for some $C_7 = C_7(d,a)>0$.
To achieve (ii), it thus suffices to adjust the choice of $\rho$ so that
\[
C_7 |\mu| (2\rho)^{p-1}\le \frac 12.
\]


With (i) and (ii) in hands, the Banach-Caccioppoli theorem guarantees the existence of a unique fixed point $U\in B_\rho$. This is the sought for mild solution to \eqref{cp02}.

\vskip 0.2in 

\noindent \underline{The $L^2_a$-mass subcritical case.} 
We next consider the subcritical case $1<p<p_c$, and let $r = p+1$. Let $q>2$ be such that $(q,r,\infty)$ is admissible, i.e.
\[
\frac 2q +\frac dr = \frac{d+a+1}2.
\]
On a finite interval $[0,T]$, with $0<T<\infty$ to be chosen subsequently, we consider the space $\XT= C([0,T],L^2_a(\RN)) \cap C^{b}_z L^q_T L^r_x$, 
with its natural norm $||\cdot||_{\XT} = ||U||_{L^\infty_T L^2_a(\RN)} + ||U||_{L^\infty_z L^q_TL^r_x }$.
In the subcritical case it is no longer true that $q =r$. In contrast with the critical case, we will show that the problem \eqref{cp02} admits a solution on $[0,T]$, regardless the size of the initial datum $u_0$ and forcing term $F$. The smallness of $T$ will enable us to close the contraction argument. 

Since $1<p<p_c = 1 + \frac{2(1-a)}{d+a+1}$, see \eqref{pcpos}, we presently must have $r<q$. To see this critical fact, note that, in view of the admissibility of $(q,r,\infty)$, $\frac 2q < \frac 2r = \frac{2}{p+1}$ is equivalent to proving
\[
\frac 2q = \frac{d+a+1}2 - \frac{d}{p+1} < \frac{2}{p+1}.
\]
Now, it is easily checked that this latter inequality is equivalent to $p < 1 + \frac{2(1-a)}{d+a+1} = p_c$, which is what we are assuming. Our next observation is that $r<q$ implies (and is in fact equivalent to)
\begin{equation}\label{pq}
p q' < q.
\end{equation}
The validity of \eqref{pq} is obvious since it is equivalent to $p<q-1\ \Longleftrightarrow\  p+1 = r <q$. At this point, we can exploit \eqref{UU}, \eqref{pq} and the smallness of the interval  
to obtain from H\"older inequality:
\begin{align}\label{Upq}
& |||U|^{p-1} U(0)||_{L^{q'}_T L^{r'}_x} 
\le T^{\delta} ||U(\cdot,0,\cdot)||^p_{L^{q}_T L^r_x}
 \le T^\delta ||U||^p_{L^\infty_z L^{q}_T L^{r}_x } \le T^\delta ||U||_{\XT}^p,
\end{align}
where we have let $\delta = \frac{1}{q'} - \frac{p}{q}>0$.
Inserting this estimate in \eqref{stelle1}, we obtain
\begin{align}\label{stelleT}
||\Lambda_a(U)||_{\XT} & \le C_3 \left[||u_0||_{L^2_a(\RN)} + ||F||_{L^1_t L^2_a(\RN)} + ||F||_{L^1_{a,z} L^{q'}_t L^{r'}_x}\right]
 +  C_1\ |\mu|\ T^\delta ||U||_{\XT}^p.
\end{align}
We now define $\rho = \rho(d,a,r, u_0, F)>0$ by the equation
\begin{equation}\label{choose}
\frac{\rho}2\ :=\ C_3 \left[||u_0||_{L^2_a(\RN)} + ||F||_{L^1_t L^2_a(\RN)} + ||F||_{L^1_{a,z} L^{q'}_t L^{r'}_x}\right],
\end{equation}
and set $B^T_\rho = \{U\in \XT\mid ||U||_{\XT}\le \rho\}$.
Inserting this choice of $\rho$ in \eqref{stelleT}, we find
\begin{align}\label{stelleT2}
||\Lambda_a(U)||_{\XT} & \le \frac{\rho}2
 +  C_1\ |\mu|\ T^\delta ||U||_{\XT}^p.
\end{align}
We next select $T = T(d,a,r,p,\mu,u_0,F)>0$ in such a  way that
\[
C_1\ |\mu|\ T^\delta \rho^p \le \frac{\rho}2.
\]
It is clear from \eqref{stelleT2} that, with such a choice of $T$, if $U\in B^T_\rho$, then $\Lambda_a(U)\in B^T_\rho$. 

Next, we show that $\Lambda_a: B^T_\rho\to B^T_\rho$ is a contraction. Proceeding as in \eqref{Larap} and \eqref{okkk}, we find
\begin{equation}\label{okkk2}
||\Lambda_a(U) - \Lambda_a(V)||_{L^\infty_{z} L^q_T L^r_x } \le C |\mu| \ |||U|^{p-1} U(0) - |V|^{p-1} V(0)||_{L^{q'}_T L^{r'}_x}.
\end{equation}
Integrating \eqref{UUU} on $[0,T]$ and using H\"older inequality, we have
\begin{align*}
& |||U|^{p-1} U(0) - |V|^{p-1} V(0)||_{L^{q'}_T L^{r'}_x} \le C(p) \left(\int_0^T \big\||U(0)| + |V(0)|\big\|^{(p-1)q'}_{L^{r}_x} ||U(0) - V(0)||^{q'}_{L^{r}_x} dt\right)^{\frac{1}{q'}}
\\
& \le C(p) \left(\int_0^T \big\||U(0)| + |V(0)|\big\|^{(p-1)q' p'}_{L^{r}_x}\right)^{\frac{1}{q'p'}} \left(\int_0^T ||U(0) - V(0)||^{q' p}_{L^{r}_x}\right)^{\frac{1}{q'p}}
\\
& \le C(p) \left(||U(0)||^{p-1}_{L^{pq'}_T L^{r}_x} + ||V(0)||^{p-1}_{L^{pq'}_T L^{r}_x}\right) ||U(0) - V(0)||_{L^{pq'}_T L^{r}_x}.
\end{align*}
In view of \eqref{pq} we obtain as in \eqref{Upq}
\begin{align*}
& \left(||U(0)||^{p-1}_{L^{pq'}_T L^{r}_x} + ||V(0)||^{p-1}_{L^{pq'}_T L^{r}_x}\right) ||U(0) - V(0)||_{L^{pq'}_T L^{r}_x} 
\\
& \le T^{\delta} \left(||U(0)||^{p-1}_{L^{q}_T L^{r}_x} + ||V(0)||^{p-1}_{L^{q}_T L^{r}_x}\right) ||U(0) - V(0)||_{L^{q}_T L^{r}_x}
\\
& \le T^{\delta} \left(||U||^{p-1}_{L^\infty_z L^{q}_T L^{r}_x} + ||V||^{p-1}_{L^\infty_z L^{q}_T L^{r}_x}\right) ||U - V||_{L^\infty_z L^{q}_T L^{r}_x}
\\
& \le T^{\delta} \left(||U||^{p-1}_{\XT} + ||V||^{p-1}_{\XT}\right) ||U - V||_{\XT}.
\end{align*}

Inserting these inequalities in \eqref{okkk2}, we finally obtain for a certain constant $C_8 = C_8(d,a,p)>0$
\begin{equation}\label{okkk3}
||\Lambda_a(U) - \Lambda_a(V)||_{L^\infty_{z} L^q_T L^r_x } \le C_8 |\mu| T^{\delta} \left(||U||^{p-1}_{\XT} + ||V||^{p-1}_{\XT}\right) ||U - V||_{\XT},
\end{equation}
and similarly
\begin{align}\label{stelle7}
||\Lambda_a(U) - \Lambda_a(V)||_{L^\infty_T L^2_a(\RN)} \le C_8\ |\mu| T^{\delta} \left(||U||^{p-1}_{\XT} + ||V||^{p-1}_{\XT}\right) ||U - V||_{\XT}. 
\end{align}
Combining \eqref{okkk3} with \eqref{stelle7}, we finally arrive to the estimate
\begin{align}\label{stelle8}
||\Lambda_a(U) - \Lambda_a(V)||_{\XT} \le 2 C_8\ |\mu| T^{\delta} \left(||U||^{p-1}_{\XT} + ||V||^{p-1}_{\XT}\right) ||U - V||_{\XT}. 
\end{align}
If we keep in mind that $\rho$ has been fixed as in \eqref{choose}, it is evident that, by possibly adjusting $T>0$ so that
\begin{equation}\label{Tdelta}
4 C_8 |\mu| \rho^{p-1} T^{\delta} \le \frac 12,
\end{equation}
we have proved that for every $U, V \in B^T_\rho$, we have
\[
||\Lambda_a(U) - \Lambda_a(V)||_{\XT} \le \frac 12 ||U-V||_{\XT}.
\]
Therefore, $\Lambda_a$ admits a unique fixed point $U\in B^T_\rho$. This proves the existence of a mild solution $U$ to the problem \eqref{cp02} in the ball $B^T_\rho$. By standard arguments based on \eqref{stelle8}, one can show that such $U$ is actually the unique mild solution in the whole space $\XT$. 

Finally, assume that $F= 0$ and that $\Im(\mu) = 0$. We want to show that, for every $M>0$, the solution $U\in \XT$ constructed above can be extended to a mild solution on $[-M,M]$, i.e.\ an element of $\mathcal X_M$ (we caution that the concatenated solution need not belong to the global space $\Xs$, since its $L^q_tL^r_x$ norm may grow with $M$). With this objective in mind, we note that \eqref{choose} shows that  $\rho=\rho(\|u_0\|_{L^2_a(\RN)})$. As a consequence of this and \eqref{Tdelta}, also $T$ depends only on $\|u_0\|_{L^2_a(\RN)}$ (besides the parameters $d, a, p$ and $\mu$). If we now solve the Cauchy problem \eqref{cp02} in $\RN\times[T,2T]$, with initial datum $U_1 = U(\cdot,\cdot,T)$, then in view of Lemma \ref{L:masscon}, we have $\|U_1\|_{{L^2_a(\RN)}}= \|u_0\|_{L^2_a(\RN)}$, and thus we can advance up to time  $2T$; that the concatenation of the two fixed points is again a mild solution on $[0,2T]$ follows from the group property of $\mathbb T^\star_a$ and the additivity of the Duhamel term in \eqref{sol1} under splitting of the time interval. The same argument applies backward in time. As a consequence, for any fixed arbitrary time interval $M>0$, with about $2M/T$ steps we obtain a mild solution on $[-M,M]$, unique in $\mathcal X_M$ by the uniqueness statement just proved on each subinterval.  

\end{proof}

\end{document}
